\input{preamble}
\usepackage{float}
\title{Detección de fauna con YOLOv8s y monitoreo de actividad en cámaras conocidas y nuevas}
\begin{document}
\maketitle
\begin{abstract}
Se ajustó YOLOv8s de Ultralytics sobre CCT20 con catorce categorías animales y
1\,000 fotografías vacías adicionales de entrenamiento. Las noventa imágenes de
la muestra de etiquetas cuentan con revisión humana. El entrenamiento terminó
 en 39 épocas y seleccionó el checkpoint de la época 29 mediante validación.
La mAP50:95 fue 0.4358 en cámaras conocidas y 0.2625 en cámaras nuevas, con una
brecha de 0.1733. Se calcularon ocho grupos de intervalos con 200 remuestreos de
secuencias cada uno y se generaron doce casos Grad-CAM, tres sin señal espacial
positiva. Un módulo propio resumió 1\,853 fotografías de una cámara nueva.
El análisis por oclusión permanece pendiente de anotación humana; no se atribuyen
resultados a niveles que todavía no fueron revisados.
\end{abstract}

\section{Motivación y trabajos relacionados}
Las cámaras trampa producen fotografías con fondos relativamente persistentes.
Recognition in Terra Incognita~\cite{terra} estudia la generalización a ubicaciones nuevas.
El objetivo es localizar animales y resumir actividad, conservando la distinción entre
 detecciones, fotografías e individuos. YOLO plantea la predicción conjunta de cajas y
clases~\cite{yolo}; se utiliza el detector YOLOv8s existente de Ultralytics~\cite{yolov8},
sin atribuir su arquitectura específica al artículo original de YOLO.
La evaluación usa AP de COCO~\cite{coco}, la optimización usa AdamW~\cite{adamw}
y la inspección exploratoria utiliza Grad-CAM~\cite{gradcam}.

\section{Benchmark, depuración y revisión humana}
Se usan las imágenes reducidas y las cinco particiones oficiales de CCT20 ECCV:
train, cis\_val, cis\_test, trans\_val y trans\_test. Las 57\,864 imágenes originales
 del benchmark no equivalen a toda la colección de Caltech. Las categorías empty y car
no son clases animales. Una imagen explícitamente vacía genera etiquetas vacías;
se excluyen vehículos sin animales y fotografías con alguna anotación animal sin caja
completa, evitando convertir objetos sin localización en negativos.

El train ECCV no contiene fotos empty. Se añadieron 1\,000 negativos del train oficial
 de la colección completa: ocho cámaras externas a las veinte del benchmark, 125
secuencias por cámara y un fotograma por secuencia. La selección usó metadatos y semilla
42, excluyó todos los identificadores y secuencias ECCV y conservó hashes y comprobaciones
 del proveedor. No se reasignaron fotografías reservadas. Estos negativos tienen
resolución original, distinta de las fotos positivas reducidas; ambos grupos se
redimensionan a 640 durante la inferencia.

Las cajas se normalizan con las dimensiones originales de metadatos. Se verifican
límites, unidades e independencia de fotografías y secuencias frente a entrenamiento;
las secuencias compartidas con cis\_val se eliminan de validación y las cámaras trans
son disjuntas de train. El conjunto elegible contiene 48\,254 imágenes. La tabla de
inventario informa fotos, fondos, cámaras y secuencias por partición.

Miguel aprobó las noventa fotografías de la muestra visual. La corrección de una
tercera ave conserva las otras dos cajas y coincide con el recuadro humano
$(681,333,707,381)$ en la referencia de $960\times700$ píxeles. Las notas individuales
 de fotos correctas son opcionales. Las decisiones, hashes y migración se conservan
sin modificar los holdouts ni los fotogramas de monitoreo.

Los originales permanecen en solo lectura. Para 1\,000 JPEG de los negativos añadidos,
V6 elimina únicamente espacios posteriores a un marcador final existente en una copia
 de trabajo. No recomprime y comprueba igualdad exacta de los píxeles decodificados.
El comprobador instalado de Ultralytics aceptó las 48\,254 imágenes sin descartar cajas.

\section{Detector y entrenamiento}
El detector principal y sus pesos preentrenados proceden de Ultralytics. La preparación,
evaluación, Grad-CAM y agregación ecológica son módulos propios. El ajuste conserva
catorce clases, resolución 640, batch 16, semilla 42, AdamW con tasa inicial
$3\times10^{-4}$, weight decay $10^{-4}$, máximo 50 épocas y paciencia 10.
Mosaic tiene probabilidad 1, el giro horizontal 0.5 y el vertical 0; se conservan
los argumentos efectivos registrados por Ultralytics.
La parada temprana se produjo en la época 39 y el mejor checkpoint fue el de la 29.
La selección utiliza cis\_val depurado y no optimiza sobre los tests.

Se conserva el desbalance observado: cinco cajas de fox y nueve de badger en train,
frente a 2\,514 de opossum. cis\_val carece de deer y fox; trans\_test carece de deer.
Una categoría sin referencia mantiene AP indefinida. Mosaic y una media global no
sustituyen la evaluación por categoría ni justifican afirmar buen desempeño en fauna rara.

\section{Protocolo de evaluación}
Se calcula mAP50 y mAP50:95 en cámaras conocidas y nuevas, con AP y soporte por
categoría. La inferencia para AP conserva confianza $\geq0.001$, NMS IoU 0.7 y máximo
300 detecciones por imagen. La brecha es $\mathrm{mAP}_{cis}-\mathrm{mAP}_{trans}$;
también se compara el mismo conjunto de trece clases con soporte en ambos tests.
El tamaño relativo usa área de caja sobre área de imagen: pequeño $<1\%$,
mediano 1--5\% y grande $>5\%$. Un sistema virtual de $1000\times1000$ preserva
IoU y permite evaluar estas áreas relativas sin confundirlas con tamaños COCO absolutos.

Los intervalos percentiles del 95\,\% usan 200 remuestreos de secuencias completas
por grupo, semilla 42 y el checkpoint fijo. Son condicionales a las cámaras observadas;
no miden variación entre entrenamientos ni incertidumbre de nuevas regiones.
Las anotaciones no proporcionan niveles fiables de oclusión. La galería preparada
contiene 17\,620 objetos elegibles de trans\_test; actualmente ninguno tiene un nivel
humano registrado. El protocolo implementado requiere revisión completa y al menos
30 objetos por nivel low, medium y high; los casos inciertos se marcan ambiguous.
Es una decisión de este protocolo, no una exigencia atribuida al docente.
Los objetos de otros niveles o ambiguos se ignoran como crowd al evaluar un grupo.
Los resultados de oclusión se mantienen pendientes hasta satisfacer esa cobertura.

\section{Resultados e interpretabilidad}
\begin{table}[H]
\centering\small
\caption{Inventario efectivo después de depurar anotaciones y secuencias. Train oficial no contiene imágenes empty. V2 añade 1.000 fondos de ocho cámaras oficiales adicionales de entrenamiento.}
\begin{tabular}{lrrrr}
\toprule
Partición & Fotos & Fondos & Cámaras & Secuencias \\
\midrule
train & 12399 & 1000 & 18 & 5733 \\
val & 2543 & 995 & 10 & 980 \\
cis\_test & 13073 & 1180 & 10 & 5378 \\
trans\_val & 1592 & 61 & 1 & 589 \\
trans\_test & 18647 & 1778 & 9 & 7438 \\
\bottomrule
\end{tabular}
\end{table}

\begin{table}[H]
\centering\small
\caption{Evaluación COCO global con el mismo checkpoint, 101 puntos de recall y máximo 300 detecciones por imagen.}
\begin{tabular}{lrrr}
\toprule
Partición & Clases con soporte & mAP50 & mAP50:95 \\
\midrule
cis\_test & 14 & 0.6878 & 0.4358 \\
trans\_test & 13 & 0.4231 & 0.2625 \\
\bottomrule
\end{tabular}
\end{table}

\begin{table}[H]
\centering\small
\caption{Rendimiento por especie: cis\_test. ND indica ausencia de referencia.}
\begin{tabular}{lrrr}
\toprule
Especie & Cajas & AP50 & AP50:95 \\
\midrule
opossum & 4008 & 0.8844 & 0.4932 \\
raccoon & 980 & 0.8970 & 0.5917 \\
squirrel & 510 & 0.7037 & 0.3696 \\
bobcat & 772 & 0.8561 & 0.6084 \\
skunk & 161 & 0.9355 & 0.6175 \\
dog & 723 & 0.6932 & 0.4858 \\
coyote & 1180 & 0.9035 & 0.6901 \\
rabbit & 1556 & 0.8874 & 0.5534 \\
bird & 752 & 0.6435 & 0.3334 \\
cat & 1528 & 0.8749 & 0.5866 \\
badger & 4 & 0.0112 & 0.0056 \\
deer & 139 & 0.5146 & 0.3383 \\
fox & 2 & 0.0000 & 0.0000 \\
rodent & 146 & 0.8236 & 0.4269 \\
\bottomrule
\end{tabular}
\end{table}

\begin{table}[H]
\centering\small
\caption{Rendimiento por especie: trans\_test. ND indica ausencia de referencia.}
\begin{tabular}{lrrr}
\toprule
Especie & Cajas & AP50 & AP50:95 \\
\midrule
opossum & 4659 & 0.8826 & 0.4590 \\
raccoon & 4769 & 0.7549 & 0.4959 \\
squirrel & 777 & 0.3644 & 0.1957 \\
bobcat & 1974 & 0.6268 & 0.4176 \\
skunk & 273 & 0.4402 & 0.2840 \\
dog & 958 & 0.7377 & 0.4994 \\
coyote & 1901 & 0.6746 & 0.4788 \\
rabbit & 722 & 0.5912 & 0.2994 \\
bird & 179 & 0.0687 & 0.0383 \\
cat & 1366 & 0.2829 & 0.1842 \\
badger & 15 & 0.0736 & 0.0583 \\
deer & 0 & ND & ND \\
fox & 1 & 0.0000 & 0.0000 \\
rodent & 26 & 0.0021 & 0.0016 \\
\bottomrule
\end{tabular}
\end{table}

\begin{table}[H]
\centering\small
\caption{Comparación con el mismo conjunto de clases con soporte en ambos tests: opossum, raccoon, squirrel, bobcat, skunk, dog, coyote, rabbit, bird, cat, badger, fox, rodent.}
\begin{tabular}{lrr}
\toprule
Partición & mAP50 & mAP50:95 \\
\midrule
cis\_test & 0.7011 & 0.4432 \\
trans\_test & 0.4231 & 0.2625 \\
\bottomrule
\end{tabular}
\end{table}

\begin{table}[H]
\centering\small
\caption{Intervalos percentiles de mAP50:95 mediante remuestreo de secuencias completas, modelo fijo. No cuantifican variación entre semillas ni entre nuevas regiones.}
\begin{tabular}{lrrr}
\toprule
Grupo & IC95 inferior & IC95 superior & Réplicas \\
\midrule
cis\_test: Global & 0.4288 & 0.4785 & 200 \\
trans\_test: Global & 0.2538 & 0.3023 & 200 \\
cis\_test: small\_lt\_1pct & 0.2452 & 0.3377 & 200 \\
cis\_test: medium\_1\_to\_5pct & 0.4516 & 0.5232 & 200 \\
cis\_test: large\_gt\_5pct & 0.5312 & 0.5720 & 200 \\
trans\_test: small\_lt\_1pct & 0.1430 & 0.1942 & 200 \\
trans\_test: medium\_1\_to\_5pct & 0.3085 & 0.3440 & 200 \\
trans\_test: large\_gt\_5pct & 0.2572 & 0.3597 & 200 \\
\bottomrule
\end{tabular}
\end{table}

\begin{table}[H]
\centering\small
\caption{Robustez por área relativa en cis\_test: menor de 1 \%, 1--5 \%, mayor de 5 \%.}
\begin{tabular}{lrrr}
\toprule
Tamaño relativo & Cajas & AP50 & AP50:95 \\
\midrule
Pequeño & 1149 & 0.5631 & 0.2792 \\
Mediano & 8925 & 0.7516 & 0.4683 \\
Grande & 2387 & 0.7583 & 0.5480 \\
\bottomrule
\end{tabular}
\end{table}

\begin{table}[H]
\centering\small
\caption{Robustez por área relativa en trans\_test: menor de 1 \%, 1--5 \%, mayor de 5 \%.}
\begin{tabular}{lrrr}
\toprule
Tamaño relativo & Cajas & AP50 & AP50:95 \\
\midrule
Pequeño & 1922 & 0.3216 & 0.1619 \\
Mediano & 12169 & 0.5272 & 0.3220 \\
Grande & 3529 & 0.4024 & 0.2870 \\
\bottomrule
\end{tabular}
\end{table}

La brecha global de mAP50:95 es de 17.33 puntos porcentuales. Con las mismas trece
categorías, la mAP pasa de 0.4432 a 0.2625: la diferencia no desaparece al igualar
la cobertura taxonómica. Las tablas conservan las clases sin soporte como ND.

\begin{figure}[H]
\centering
\includegraphics[width=0.74\linewidth]{figures/figure_0.png}
\includegraphics[width=\linewidth]{figures/figure_1.png}
\caption{Distribución de cajas de entrenamiento y evolución de pérdidas y mAP en cis\_val. El mejor checkpoint corresponde a la época 29; las curvas no utilizan los tests.}
\end{figure}

\begin{figure}[H]
\centering
\includegraphics[width=\linewidth]{figures/figure_2.png}
\caption{Seis casos en cámaras conocidas. Se conserva el mapa plano de cis\_test\_3 como limitación del método; no se interpreta como evidencia de ausencia de atención. Los rótulos day/night son aproximaciones de metadatos.}
\end{figure}

\begin{figure}[H]
\centering
\includegraphics[width=\linewidth]{figures/figure_3.png}
\caption{Seis casos en cámaras nuevas. trans\_test\_3 y trans\_test\_5 tienen mapas planos. Las zonas resaltadas en fondos y bordes no demuestran causalidad ni memorización.}
\end{figure}

Grad-CAM fija una propuesta y el score diferenciable de una clase antes de NMS,
con una capa del backbone y el letterbox declarados. Hay seis casos por test que
cubren aciertos a IoU $\geq0.5$, omisiones o clase incorrecta y falsos positivos en
fondos. La selección prioriza categorías raras y frecuentes mediante reglas
 deterministas. Tres mapas tienen rango $[0,0]$ y no ofrecen una explicación espacial.
Los rótulos day/night proceden de la hora en metadatos; varias fotos rotuladas day
son oscuras y de apariencia infrarroja. No se infiere iluminación real de ese rótulo
ni se presenta una comparación cuantitativa por iluminación.

\section{Demostración de monitoreo}
La cámara trans 7 aporta todos sus 1\,853 fotogramas oficiales, incluidos los excluidos
 de la evaluación de cajas. Con confianza $\geq0.25$, 1\,420 fotografías contienen
 detecciones y la tasa predicha de fotos vacías es 23.37\,\%. Se resumen 211 fechas
con fotografías y once categorías detectadas. Se conserva la zona horaria sin
especificar de los metadatos; no se inventan fechas ni tiempo de exposición.
La ráfaga de demostración tiene tres fotogramas consecutivos del 8 de mayo de 2011:
los dos primeros no tienen detecciones y el tercero registra opossum.
\begin{table}[H]
\centering\small
\caption{Actividad predicha de la cámara 7: primeros diez días de la cohorte completa de esa cámara.}
\begin{tabular}{lrrrr}
\toprule
Día & Fotos & Con fauna & Especies & Vacías \\
\midrule
2011-05-07 & 1 & 0 & 0 & 1.0000 \\
2011-05-08 & 9 & 2 & 1 & 0.7778 \\
2011-05-10 & 3 & 1 & 1 & 0.6667 \\
2011-05-20 & 3 & 3 & 1 & 0.0000 \\
2011-05-25 & 2 & 1 & 1 & 0.5000 \\
2011-06-03 & 3 & 1 & 2 & 0.6667 \\
2011-06-05 & 3 & 3 & 1 & 0.0000 \\
2011-06-17 & 2 & 1 & 1 & 0.5000 \\
2011-06-19 & 6 & 3 & 2 & 0.5000 \\
2011-06-20 & 15 & 14 & 2 & 0.0667 \\
\bottomrule
\end{tabular}
\end{table}

\begin{table}[H]
\centering\small
\caption{Resumen de todos los días de la cámara de demostración. Las detecciones no representan individuos únicos.}
\begin{tabular}{lrr}
\toprule
Especie & Detecciones & Fotos con especie \\
\midrule
bird & 92 & 91 \\
bobcat & 49 & 49 \\
cat & 45 & 43 \\
coyote & 477 & 449 \\
deer & 1 & 1 \\
dog & 33 & 32 \\
opossum & 649 & 643 \\
rabbit & 81 & 80 \\
raccoon & 90 & 85 \\
skunk & 10 & 10 \\
squirrel & 38 & 34 \\
\bottomrule
\end{tabular}
\end{table}

\begin{table}[H]
\centering\small
\caption{Secuencia completa con fecha de captura válida. Identificadores completos disponibles en activity\_sequence.csv.}
\begin{tabular}{lrrr}
\toprule
Fotograma & Captura & Detecciones & Especies \\
\midrule
1 & 2011-05-08 11:39:22 & 0 & ND \\
2 & 2011-05-08 11:39:23 & 0 & ND \\
3 & 2011-05-08 11:39:24 & 1 & opossum \\
\bottomrule
\end{tabular}
\end{table}

Una foto con varias cajas de la misma categoría se cuenta una vez en fotos con
categoría. bird y rodent agrupan taxones, por lo que el conteo de categorías no mide
 diversidad biológica completa. Las detecciones no representan individuos únicos y
la tasa predicha de vacíos no equivale a una revisión humana de ausencia.

\section{Discusión y uso responsable}
La mAP50:95 disminuye de 0.4358 en cámaras conocidas a 0.2625 en cámaras nuevas.
Los intervalos del 95\,\% son $[0.4288,0.4785]$ y $[0.2538,0.3023]$,
respectivamente. La comparación con trece categorías comunes conserva una brecha
de 18.08 puntos porcentuales. Es consistente con dificultad de transferencia
entre cámaras, pero los intervalos no son una prueba causal del efecto del fondo
ni estiman variación entre semillas o nuevas regiones. El modelo fue seleccionado
en validación; los resultados de prueba se conservan como evaluación de ese ajuste.

Los objetos pequeños tienen mAP50:95 de 0.2792 en cis\_test y 0.1619 en trans\_test.
En cámaras nuevas, los medianos alcanzan 0.3220 y los grandes 0.2870; el rendimiento
no crece monótonamente con el área. La composición de categorías y escenas difiere
entre estratos, por lo que el tamaño no se aísla como una causa mediante estas tablas.
bird pasa de AP50:95 0.3334 a 0.0383 y rodent de 0.4269 a 0.0016, con 179 y 26
cajas de prueba trans, respectivamente. Estas caídas necesitan evaluación específica.
Fox obtiene AP cero con solo dos referencias cis y una trans, después de cinco cajas
de entrenamiento; badger tiene nueve cajas de entrenamiento. Esos soportes impiden
considerar la media global como evidencia de detección fiable de categorías raras.

Los doce casos Grad-CAM incluyen nueve mapas con variación espacial y tres planos.
En algunos aciertos se resaltan partes del animal; los falsos positivos también
resaltan vegetación, suelo o bordes de la entrada. Estas observaciones son exploratorias.
Los mapas planos de opossum no permiten localizar contribuciones positivas para
la propuesta y capa fijadas. No se reemplazaron los casos ni se modificó la técnica
congelada para obtener visualizaciones favorables. Los rótulos de horario no son
mediciones de iluminación, especialmente cuando una fotografía oscura se rotula day.
La galería no estima la frecuencia de estos patrones ni demuestra dependencia causal
del fondo. Para comprobarla harían falta intervenciones controladas sobre las imágenes.

La oclusión queda sin resultado: los 17\,620 objetos elegibles no tienen niveles
humanos registrados. No se deduce oclusión a partir del área, la confianza o el mapa.
Los análisis globales y por tamaño están completos, pero no sustituyen esa revisión.
La demostración de la cámara 7 muestra una agregación reproducible de fotografías;
sus 1\,420 fotos con detecciones y once categorías no son un censo de individuos.

La ampliación v2 incorpora solo negativos en ocho cámaras adicionales; no agrega animales de esas cámaras.
El fondo puede entonces ayudar a distinguir la procedencia de los negativos. La resolución original
de esa cohorte también difiere de las imágenes ECCV reducidas. El resultado se describe como
entrenamiento ampliado y no como reproducción estricta del train ECCV. Una comparación futura
necesitaría otra ejecución controlada, conservando los mismos tests y sin ajustar a sus resultados.
Un falso negativo de una especie amenazada puede invisibilizar su presencia en decisiones de manejo,
mientras que un falso positivo consume tiempo de revisión y puede crear registros de distribución incorrectos.
El umbral 0.25 de la demostración no se considera óptimo para todas las especies ni para una decisión de conservación.

Debe mantenerse revisión profesional de observaciones sensibles, documentar esfuerzo de muestreo
y validar en cámaras y hábitats del destino. Una tasa baja de detecciones también puede reflejar
fallos del sensor o cambio de iluminación. Los conteos de imágenes no estiman población ni individuos.
La separación entre cámaras ayuda a evaluar generalización, pero el entrenamiento en Estados Unidos
no justifica despliegue en Amazonía o Andes sin validación local.

La decisión de revisar una fotografía debe considerar el costo de omitir una observación,
el costo de una identificación incorrecta y la capacidad del equipo de campo para verificarla.
Un umbral global puede facilitar una demostración operativa, pero no sustituye un protocolo
por especie. Para registros sensibles conviene conservar la fotografía original, la caja,
la confianza y la decisión del experto, evitando que una etiqueta automática se convierta
sin revisión en un registro confirmado. El procesamiento debe poder repetirse si cambian
el modelo o el criterio de revisión.

La concentración de ejemplos en especies frecuentes puede ocultar un rendimiento insuficiente
en especies raras aun cuando la mAP global resulte favorable. Los soportes reales, las categorías
sin referencia y las exclusiones deben acompañar cada comparación; una AP indefinida no es
evidencia de ausencia biológica. La revisión de casos nocturnos, objetos pequeños y cuerpos
parcialmente visibles ayuda a priorizar nuevas muestras, pero una galería exploratoria no mide
por sí sola la frecuencia de esos errores. Para una aplicación de conservación se necesitarían
muestreo temporal representativo, control del esfuerzo por cámara y evaluación en el hábitat
de destino. Cualquier decisión de manejo conserva responsabilidad humana y debe admitir
corrección cuando aparezcan evidencias adicionales.


\section{Reproducibilidad y pendientes}
Se conservaron 44 fuentes experimentales congeladas, manifiestos, exclusiones,
argumentos efectivos, pesos, predicciones y métricas. Las dependencias están fijadas
 y se verificaron CUDA, NMS, gradientes y AMP en RTX 5090. La prueba pequeña de
reanudación comprueba época, optimizador y paciencia; no afirma equivalencia bit a bit.
El entrenamiento científico y esa prueba tienen salidas independientes.

Un límite de memoria interrumpió el cálculo inicial de intervalos después del
entrenamiento. La recuperación ejecutó cada grupo en un proceso nuevo con la función
original, conservando semilla, réplicas, datos y modelo. No repitió el entrenamiento;
las pruebas verificaron igualdad exacta de resultados en casos de referencia.
Los ocho grupos quedaron guardados. El notebook ejecutó sus doce celdas de código
sin errores y la demostración y los doce casos Grad-CAM están disponibles.
El código está en \url{https://github.com/miguelsff/maestria-ia-uni-ciclo-2-redes-neuronales}.
Este documento conserva carácter de borrador: faltan niveles humanos de oclusión
 y aprobación editorial del equipo. No se registra aprobación del docente.

\bibliographystyle{plain}
\bibliography{references}
\end{document}
